\section{Endpoint degrees and small diagonal support}\label{sec:endpoints}

\subsection{Chow transfer and complementary degrees}

\begin{lemma}[Transfer through a Chow diagonal]\label{lem:general-transfer}
Suppose that over a finite extension $L/K$ there are smooth projective
varieties $V_j$ and correspondences of opposite degrees satisfying
\[
 \Delta_{X_L}=\sum_jv_j\circ u_j
 \quad\text{in }\CH^{\dim X}(X_L\times_LX_L)_\Q.
\]
Fix $r$. If every homologically trivial image $u_{j*}z_L$ is in the
native obstruction kernel of $V_j$, then
$z\in\cK_K^r(X)$. In particular $\mathsf V(V_j)$ for every $j$
implies $\mathsf V(X)$.
\end{lemma}

\begin{proof}
Cycle classes commute with correspondences, so the image classes are
homologically trivial. Kernel covariance sends $v_{j*}u_{j*}z_L$ into
$\cK_K^r(X_L)$. Sum the diagonal and push down by $X_L\to X$.
The result is $[L:K]z$, and rational division gives the conclusion.
All correspondences over $L$ are interpreted on the diagonal
open-and-closed factor of the product over $K$, as in
\subref{sub:covariance}. Finally apply the curved-pair theorem to the
specified representative of $z$.
\end{proof}

\begin{lemma}[The two-degree criterion]\label{lem:two-degrees}
Let $X/F$ have dimension $n$, and let $S\subseteq\{1,\ldots,n\}$.
Suppose that at every finite place $v$, all geometric homological
classes in degrees $r\in S$ have zero native obstruction. Then
$\cH^r(X)=\CH^r(X)_{\Q,\geomhom}$ for $r\in S$. The full height
pairing in degrees $(r,n+1-r)$ is therefore defined whenever
$r\in S\cap(n+1-S)$.
\end{lemma}

\begin{proof}
Geometric cycle classes commute with passage to $F_v$. Thus every
global homologically trivial class in an asserted degree belongs to
each local kernel defining $\cH^r(X)$. The opposite inclusion is in
the definition of $\cH^r(X)$. Apply the companion pairing to both
resulting equalities.
\end{proof}

\begin{proposition}[Endpoint input]\label{prop:endpoint-input}
For a smooth projective $n$-fold over the original $K$, geometric
homological divisors, zero-cycles and one-cycles have zero obstruction.
Thus $\mathsf V(X)$ holds if $n\leq3$.
\end{proposition}

The divisor and zero-cycle assertions are the endpoint results in
\cite{InternalEndpoints}. For one-cycles, the precise input is the
regular-model curve correction followed by dominant obstruction
descent, also recorded there. Its proof proceeds as follows. Choose a
projective strictly semistable alteration $V\to X$, over a finite
constant extension, of the same dimension as $X$. The pullback of a
homologically trivial one-cycle remains a one-cycle. The curve
correction produces a rational virtual bundle on the regular model
whose generic degree-$(n-1)$ character is that class and whose degree
on every vertical $(n-1)$-cycle is zero. The numerical-model criterion
annihilates its obstruction; dominant descent gives the assertion on
$X$. This uses the full degree of the alteration and leaves the final
curved-pair correction on $X$.

\subsection{Two orientations of the diagonal}

\begin{definition}\label{def:small-support}
A rational small-support diagonal for $X$, of dimension $n$, is an
identity
\begin{equation}\label{eq:small-diagonal}
 \Delta_X=A+B,\qquad
 \operatorname{Supp}A\subseteq D\times X,\quad
 \operatorname{Supp}B\subseteq X\times T,
\end{equation}
in rational Chow, where $D$ is a divisor and $\dim T\leq3$.
The identity may be supplied after a finite extension of the ground
field. The supports belong to actual cycle representatives of $A,B$.
\end{definition}

\begin{theorem}[Small diagonal support]\label{thm:small-support}
If $X$ has a small-support diagonal, its complete geometric homological
Chow groups are in the local obstruction kernel for
\begin{equation}\label{eq:small-support-range}
 r\leq2\quad\text{or}\quad r\geq n-2.
\end{equation}
It has property $\mathsf V$ if $n\leq5$. Over a number field, the
same geometric hypothesis gives a full height domain if $n\leq5$;
for $n\geq6$ it gives the full pairs with $r\leq2$ or $r\geq n-1$.
\end{theorem}

\begin{proof}
Resolve the components of $D$ and $T$, with maps $a_i:D_i\to X$ and
$b_j:T_j\to X$. A cycle supported on the image of a proper surjection
can be lifted rationally: choose a closed point of the fiber above its
generic point, take its closure, and divide by the full generic
degree. This also applies to cycles contained in exceptional images.
It lifts \eqref{eq:small-diagonal} to correspondences
\[
 U_i\in\CH^{n-1}(D_i\times X)_\Q,\qquad
 V_j\in\CH^{t_j}(X\times T_j)_\Q,\quad t_j=\dim T_j\leq3.
\]
The identity and its transpose act as
\begin{align}
 z&=\sum_iU_{i*}a_i^*z+\sum_jb_{j*}V_{j*}z,
       \label{eq:small-forward}\\
 z&=\sum_i a_{i*}{}^tU_{i*}z+\sum_j{}^tV_{j*}b_j^*z.
       \label{eq:small-transpose}
\end{align}
All intermediate classes are geometrically homologically trivial.
In \eqref{eq:small-forward}, the class on $D_i$ has dimension
$n-1-r\leq1$ when $r\geq n-2$. In
\eqref{eq:small-transpose}, it has codimension $r-1\leq1$ when
$r\leq2$. These are the endpoint degrees. Every feasible class on
$T_j$ is an endpoint class because $t_j\leq3$; impossible degrees
and homologically trivial rank classes are zero. Apply kernel
covariance to the indicated orientation. If the identity is defined
after a finite extension, use \lemref{lem:general-transfer} to descend.
The two local ranges cover every degree for $n\leq5$. Intersecting
them with their complementary ranges proves the height assertions
by \lemref{lem:two-degrees}.
\end{proof}

\begin{proposition}[Universal zero-cycle support]\label{prop:universal-support}
Suppose that a closed subset $T\subset X_{\bar k}$, of dimension at
most $s$, surjects onto $\CH_0(X_\Omega)_\Q$ for every algebraically
closed extension $\Omega/\bar k$. Then after a finite extension of
$k$ there is an identity \eqref{eq:small-diagonal} with $\dim T\leq s$.
\end{proposition}

\begin{proof}
Apply the support hypothesis to the generic point over
$\overline{\bar k(X)}$. A rational-equivalence certificate uses only
finitely many points, curves and functions, so it is defined over a
finite extension of $\bar k(X)$. Norming down and dividing by its
degree gives the supported generic-point relation over $\bar k(X)$.
Close the cycle in $X_{\bar k}\times T$ and spread the relation over
an open subset of the first factor. Chow localization puts the
remaining difference from $\Delta_X$ over its closed complement;
contain this complement in a divisor. Finally descend the finite
cycle-and-function certificate and the supports to a finite
extension of $k$.
\end{proof}

\begin{corollary}[Rationally connected fibrations]\label{cor:rc-fibrations}
Let $f:X\to B$ be a surjective morphism of smooth projective varieties
in characteristic zero. Suppose that over a nonempty open of $B$ it
is smooth with rationally connected geometric fibers. If $\dim B\leq3$,
then \thmref{thm:small-support} applies to $X$. In particular it gives
all degrees if $\dim X\leq5$. Geometrically rationally connected
varieties, including smooth Fano varieties, are examples with
zero-dimensional support.
\end{corollary}

\begin{proof}
The closure of a closed point of the generic fiber is a multisection
$T\subset X$ of dimension $\dim B$. It is proper and surjective over
$B$. After any algebraically closed extension, move a zero-cycle
into the smooth rationally connected locus. Each of its points is
rationally equivalent, inside its fiber, to a point of $T$ in the
same fiber. Rational connectedness gives this equality by chains of
rational curves. The multisection therefore supplies universal
geometric rational zero-cycle support. Apply
\propref{prop:universal-support} and \thmref{thm:small-support}.
\end{proof}

\subsection{Projective bundles and dominant images}

\begin{proposition}\label{prop:closures}
Property $\mathsf V$ is preserved by projective bundles, by blowups
along smooth centers when both the ambient variety and the center
have $\mathsf V$, and by smooth projective dominant morphism images.
The assertions allow a preliminary finite extension.
\end{proposition}

\begin{proof}
The Chow-motive projective-bundle and blowup decompositions are
finite direct sums of Tate twists of the indicated smooth varieties.
Apply \lemref{lem:general-transfer}. For a projective dominant
morphism $f:V\to X$ with relative dimension $e$, choose a relatively
ample class $a$ whose top degree on the generic fiber is $d>0$.
The correspondence $d^{-1}f_*(a^e\cdot-)$ is a left inverse to
$f^*$ by the projection formula. Its action gives a rational Chow
retraction and hence the same transfer. The formula can be checked
after product with any smooth projective test variety, so it is a
Chow correspondence identity. Finite descent is as before.
\end{proof}

\begin{corollary}\label{cor:uniruled}
A geometrically uniruled smooth projective fourfold has $\mathsf V$.
So does any smooth projective dominant morphism image of a
projective-bundle tower over a variety of dimension at most three,
or over one of these uniruled fourfolds.
\end{corollary}

\begin{proof}
After a finite extension, a uniruled fourfold is dominated rationally
by $B\times\PP^1$ for a smooth projective threefold $B$. Resolve this
map in characteristic zero by smooth blowups of the four-dimensional
source. Every nontrivial center has dimension at most two, hence has
$\mathsf V$. The source and its resolution have $\mathsf V$ by
\propref{prop:closures}. The resolved dominant morphism transfers it
to $X$. The remaining assertions follow successively from projective
bundle and dominant-morphism transfer.
\end{proof}
